\documentclass[a4paper]{article}

% Español (México) con XeLaTeX: fontspec para fuentes Unicode, polyglossia
% para la división silábica y los nombres del documento en la variante mexicana.
\usepackage{fontspec}
\usepackage{polyglossia}
\setdefaultlanguage[variant=mexican]{spanish}
% Los nombres chinos van como «pinyin (original)»: xeCJK da a los caracteres
% Han su propia fuente; sin ella XeLaTeX los omite con un aviso.
% FandolSong no cubre algunos caracteres de nombres (U+73FA); WenQuanYi Zen Hei sí.
\usepackage[AutoFallBack=true]{xeCJK}
\setCJKmainfont{FandolSong-Regular.otf}
\setCJKfallbackfamilyfont{\CJKrmdefault}{WenQuanYi Zen Hei}
% polyglossia no traduce los nombres de listings.
\AtBeginDocument{\providecommand{\lstlistingname}{}\renewcommand{\lstlistingname}{Listado}%
  \providecommand{\lstlistlistingname}{}\renewcommand{\lstlistlistingname}{Índice de listados}}
\usepackage{amsmath,amssymb}
\usepackage{graphicx}
\usepackage[margin=2.4cm]{geometry}
\usepackage[most]{tcolorbox}
\usepackage{etoolbox}
\usepackage{subcaption}
\usepackage{listings}
\usepackage[hidelinks]{hyperref}
\usepackage{booktabs}
\usepackage{float}
\usepackage{enumitem}
\usepackage{tabularx}
\usepackage{caption}
\setlist[itemize]{topsep=0.35em, itemsep=0.25em}
\captionsetup{font=small,labelfont=bf}

\newtcolorbox{knowledgebox}[1]{
    enhanced, colback=blue!5!white, colframe=blue!75!black, colbacktitle=blue!75!black,
    coltitle=white, fonttitle=\bfseries, title={#1}, boxrule=1pt, sharp corners
}
\newtcolorbox{importantbox}[1]{
    enhanced, colback=yellow!10!white, colframe=yellow!80!black, colbacktitle=yellow!80!black,
    coltitle=black, fonttitle=\bfseries, title={#1}, sharp corners, boxrule=1pt
}
\newtcolorbox{warningbox}[1]{
    enhanced, colback=red!5!white, colframe=red!75!black, colbacktitle=red!75!black,
    coltitle=white, fonttitle=\bfseries, title={#1}, sharp corners, boxrule=1pt
}

\lstset{
    basicstyle=\ttfamily\small,
    keywordstyle=\color{blue},
    stringstyle=\color{red!60!black},
    commentstyle=\color{green!60!black},
    breaklines=true,
    frame=single,
    numbers=left,
    numberstyle=\tiny\color{gray},
    captionpos=b,
    extendedchars=true
}

\newcommand{\notetitle}{CS224R Lecture 16: Autonomy — Chelsea Finn}
\newcommand{\noteauthors}{Elaborado a partir de los materiales públicos de Stanford CS224R}
\newcommand{\notedate}{18 de marzo de 2025}
\newcommand{\videochannel}{Stanford CS224R: Reinforcement Learning}
\newcommand{\videopublishdate}{18 de marzo de 2025}
\newcommand{\videoduration}{aproximadamente 80 minutos}
\newcommand{\videourl}{https://cs224r.stanford.edu/}
\newcommand{\videocoverpath}{cover.jpg}
\newcommand{\repourl}{https://github.com/hqhq1025/ai-course-notes}

\usepackage{fancyhdr}
\pagestyle{fancy}
\fancyhf{}
\fancyfoot[C]{\thepage}
\fancyfoot[R]{\footnotesize\color{gray}\href{\repourl}{AI Course Notes}}
\renewcommand{\headrulewidth}{0pt}
\renewcommand{\footrulewidth}{0pt}

\begin{document}
\begin{titlepage}
\centering
{\Large Notas del curso\par}
\vspace{1.2cm}
{\huge\bfseries \notetitle\par}
\vspace{0.8cm}
{\large \noteauthors\par}
\vspace{0.3cm}
{\large \notedate\par}
\vspace{1.1cm}

\ifdefempty{\videocoverpath}{
{\small No se proporcionó imagen de portada.\par}
}{
\includegraphics[width=0.82\textwidth,height=0.45\textheight,keepaspectratio]{\videocoverpath}\par
}

\vfill
\begin{tcolorbox}[width=0.9\textwidth, colback=black!2!white, colframe=black!60, sharp corners]
\textbf{Autor o canal del video}: \videochannel\par
\textbf{Fecha de publicación}: \videopublishdate\par
\textbf{Duración del video}: \videoduration\par
\textbf{Ponente}: Chelsea Finn (investigadora de Stanford CS224R)\par
\textbf{Enlace del video}: \href{\videourl}{\nolinkurl{\videourl}}\par
\textbf{Página del proyecto}: \href{\repourl}{\nolinkurl{\repourl}}\par
\end{tcolorbox}
\end{titlepage}

\tableofcontents
\newpage
\section{La lógica del aprendizaje robótico autónomo}
\subsection{La brecha entre Human-in-the-loop y Autonomy}
El sistema actual de aprendizaje robótico depende de la supervisión humana: la función de recompensa, el reinicio del entorno y la recolección de datos están a cargo de ingenieros. La clase «Autonomy» plantea un reto: que el robot explore activamente el complejo mundo físico como lo hace un bebé, en lugar de esperar a que una persona le fije los objetivos.

\begin{figure}[H]
\centering
\includegraphics[width=0.92\textwidth]{slides-images/slide-01.jpg}
\caption{La Slide 1 presenta el tema central de la clase: Autonomy es, a la vez, el conjunto de la configuración de reward, reset y goal.}
\end{figure}

\begin{importantbox}{Los tres objetivos de Autonomy}
1. Construir la recompensa de forma automática o entender el resultado de la tarea;\\
2. Operar de forma continua en un mundo sin reinicios;\\
3. Elegir objetivos de forma autónoma y explorar con seguridad.
\end{importantbox}

\subsection{Criterios de evaluación frente a la realidad}
El ponente enfatiza tres dimensiones de evaluación: \textbf{capacidad de aprendizaje} (si se puede obtener suficiente signal), \textbf{escalabilidad} (si no depende de reinicios manuales) y \textbf{sostenibilidad} (si puede funcionar a largo plazo y transferirse a objetivos nuevos). La lógica de la exposición: se parte de la reward y, después, a partir de la lógica del entorno y de los objetivos, se avanza hacia la seguridad y la integración de sistemas.

\begin{knowledgebox}{Mapa de la lógica de la exposición}
La línea de investigación de Autonomy sigue este orden: \textbf{recompensa automática}→\textbf{control Reset-free}→\textbf{Goal Management}→\textbf{Safe Exploration}→\textbf{monitoreo del sistema}. Cada etapa parte de un supuesto anterior (reward, reset, goal) y describe cómo retirarlo.
\end{knowledgebox}

\subsection{Resumen de la sección}
Esta sección repasa la situación actual y la lógica de la exposición del aprendizaje robótico autónomo, y toma las cinco etapas «recompensa-entorno-objetivo-seguridad-sistema» como marco rector de las secciones siguientes.
\newpage
\section{Recompensas automáticas y señales de comprensión}
\subsection{Representación del objetivo con imágenes y lenguaje}
El requisito central de las recompensas autónomas es reducir el trabajo de ingeniería manual. Chelsea describió el pipeline para automatizar la recompensa con una \textbf{imagen objetivo} o una \textbf{descripción en lenguaje}: un modelo preentrenado extrae la semántica visual, se compara el escalamiento de los embedding del estado actual y del objetivo, y así se obtiene una recompensa de tipo arrastre.

\begin{figure}[H]
\centering
\includegraphics[width=0.9\textwidth]{slides-images/slide-02.jpg}
\caption{La Slide 2 muestra cómo la distancia entre los embedding del objetivo y del estado actual funciona como reward inmediata.}
\end{figure}

\begin{importantbox}{Embedding Distance Reward}
\begin{itemize}
\item Se usan modelos como VAE/CLIP para extraer `f(I)`;\\
\item La recompensa es $r_t = -\|\phi(I_t) - \phi(I_g)\|_2$;\\
\item Las diferencias de dominio pueden atenuarse con scratch/finetune, pero lo esencial es comparar estados mediante la representation.
\end{itemize}
\end{importantbox}

\subsection{Combinación de reward visual y de lenguaje}
La Slide 03 muestra el flujo que combina descripciones en lenguaje natural con embedding visuales para generar la reward: el LLM genera el task prompt, el VLM evalúa el current frame con base en el prompt y, al final, se filtra con un ensemble output. Con unos cuantos cases especiales de human-in-the-loop, el sistema puede distinguir las reward signal inseguras.

\begin{figure}[H]
\centering
\includegraphics[width=0.85\textwidth]{slides-images/slide-03.jpg}
\caption{La Slide 3 muestra el sistema en el que un LLM y un VLM se combinan para producir la reward, y subraya la importancia de la etapa de filter.}
\end{figure}

\begin{knowledgebox}{LLM + VLM Reward Pipeline}
\begin{itemize}
\item El LLM recibe el task prompt y genera una structured description;\\
\item El VLM (por ejemplo, CLIP/BLIP) califica el frame actual;\\
\item Un confidence filter (por ejemplo, la softmax temperature) controla la intensidad de la reward.
\end{itemize}
\end{knowledgebox}

\subsection{Análisis detallado del Reward Pipeline}
A partir del flujo de la slide 03, el expositor detalló las tres etapas del reward pipeline: \textbf{signal extraction} (obtener imágenes/lenguaje de los sensors), \textbf{affinity scoring} (embedding similarity) y \textbf{confidence gating} (filtrado por confianza y ensemble). Estas tres etapas deben modificarse de forma sincronizada después de cada actualización del modelo para garantizar que la reward not drift.

\begin{importantbox}{Marco abstracto del pipeline}
\begin{itemize}
\item Signal extraction: logging sensors + normalization;\\
\item Affinity scoring: distancia entre embedding o classifier probability;\\
\item Confidence gating: usar un filtro de threshold/ensemble/temporal smoothing para descartar las reward anómalas.
\end{itemize}
\end{importantbox}

\subsection{Autosupervisión basada en clasificadores de éxito}
El expositor recurre a anotaciones few-shot para que el robot aprenda a reconocer estados de éxito: se reúnen unas cuantas instantáneas positive/negative, se entrena un clasificador binario $\sigma(s)$ y la salida de ese classifier es la función de reward. Este enfoque es muy adecuado para tareas en las que «se desconoce el número de objetivos».

\begin{knowledgebox}{Success Classifier Pipeline}
\begin{itemize}
\item Reunir de 50 a 100 frames de «éxito» definidos por el usuario y generar muestras negativas;\\
\item Entrenar una CNN/RNN ligera que produzca probabilidades;\\
\item Usar la probabilidad directamente como reward, o combinarla con un decaimiento temporal para evitar una convergencia prematura.
\end{itemize}
\end{knowledgebox}

\subsection{Filtrado de reward multimodal}
Chelsea señaló además que, cuando CLIP/LLM generan la reward, es indispensable añadir un \textbf{filtrado por confiabilidad} para evitar que el modelo asigne puntajes altos a irrelevant features (como el fondo). Las prácticas habituales son hacer un ensemble de varias modalidades y agregar una regularización de temporal consistency.

\begin{warningbox}{No dejes que los resultados de hallucination del modelo pretrained lleguen a la policy}
Una reward basada en CLIP puede reforzar destellos del fondo. Se recomienda:
\begin{itemize}
\item Usar temporal smoothing para asegurar que la reward sea continua;\\
\item En los juicios multimodales, emplear majority vote/variance penalty;\\
\item Conservar un canal de audit human-in-the-loop para hacer override en caso de emergencia.
\end{itemize}
\end{warningbox}

\subsection{Resumen de la sección}
Las recompensas automáticas sustituyen la reward manual tradicional con diversas señales, como imágenes, lenguaje y clasificadores basados en modelos, y subrayan la necesidad de controlar la confianza de las salidas multimodales.
\newpage
\section{Control reset-free y gestión de la memoria}
\subsection{Debilitamiento del supuesto de reinicio}
El RL tradicional depende de reinicios hechos por una persona o por el simulator. Chelsea propone «que la propia policy haga el reset»: al terminar cada episodio de la tarea se activa otra policy que devuelve el entorno a un nuevo conjunto inicial, de modo que el robot pueda operar de forma continua.

\begin{figure}[H]
\centering
\includegraphics[width=0.88\textwidth]{slides-images/slide-04.jpg}
\caption{La slide 4 explica cómo los controladores hacia adelante y hacia atrás alternan entre ejecutar la tarea y hacer el reset.}
\end{figure}

\noindent La timeline de la figura representa el proceso de recuperación después de una ejecución de la tarea: cuando la policy de la tarea (al centro) alcanza el goal, entrega el current state a la reset policy. Al ejecutarse, la reset policy lleva el system de vuelta a la base manifold y sigue proporcionando puntos de partida para la tarea.

\begin{importantbox}{Mecanismo del Forward-Backward Controller}
\begin{itemize}
\item $\pi_f$: explora, a partir del initial state, una task-specific trajectory;\\
\item $\pi_b$: recibe el state actual e intenta volver a la initial manifold;\\
\item Se entrenan ambas por turnos; la tasa de éxito de $\pi_b$ define el punto de partida del nuevo episode.
\end{itemize}
\end{importantbox}

\subsection{Episodic buffer persistente}
El enfoque reset-free requiere un long-lived buffer que recuerde los «failed states más recientes», para que la política siga aprendiendo a hacer el reinicio. Chelsea recomienda mantener un prioritized replay buffer e incorporar además \textbf{time-aware sampling}, que da prioridad a reproducir los datos generados después del reset más reciente.

\begin{knowledgebox}{Time-Aware Replay}
\begin{itemize}
\item Se añade un timestamp a cada transition;\\
\item Durante el entrenamiento se usa $w_t \propto \exp(-\lambda (t_{now}-t))$ para aproximar la latest priority;\\
\item Así se garantiza que los datos nuevos se aprovechen plenamente, sin dejar de conservar algunos rare events.
\end{itemize}
\end{knowledgebox}

\subsection{Deshacer experiencia y diversidad del entorno}
El ponente señala además que un sistema autónomo debe tener la capacidad de \textbf{deshacer experiencia}, es decir, que el modelo retroceda sobre las trayectorias en las que hizo fail y pruebe nuevas combinaciones de estrategias. Esto implica planificar de forma dinámica varias reset-policies para cubrir distintas perturbaciones físicas.
\begin{warningbox}{Fail-fast but recover-fast}
Se alienta a la política a probar con rapidez varias combinaciones de acciones y, si el control actual falla, a activar de inmediato la reset policy; al mismo tiempo, se registran los modos de falla para la future policy adaptation.
\end{warningbox}

\subsection{Resumen de la sección}
El marco reset-free usa forward-backward controllers, un time-aware buffer y fast recovery para lograr una operación continua sin depender de resets manuales.
\newpage
\section{Establecimiento autónomo de objetivos y aprendizaje por currículo}
\subsection{Banco de objetivos y curriculum jerárquico}
Chelsea utiliza un «goal replay buffer» para registrar los objetivos alcanzados o casi alcanzados y, junto con una \textbf{curriculum policy}, selecciona para el entrenamiento objetivos de dificultad moderada. Cada vez que la policy alcanza un objetivo, este se escribe, junto con sus indicadores de capacidad, en el «tracker de estado de objetivos».

\begin{figure}[H]
\centering
\includegraphics[width=0.88\textwidth]{slides-images/slide-05.jpg}
\caption{La Slide 5 explica cómo implementar un curriculum automático en goal-conditioned RL.}
\end{figure}

\begin{tabularx}{\textwidth}{lX}
\toprule
Etapa & Medida principal \\
\midrule
Exploración & Al muestrear objetivos de entrenamiento, se da preferencia a novel/uncertain states; \\
Equilibrio & Se fija una success rate window (por ejemplo, 50\%-80\%) para mantener la zona de aprendizaje; \\
Avance & Por cada objetivo logrado, se aplica un parameterized difficulty increment; \\
\bottomrule
\end{tabularx}

\subsection{Goal-Query y el razonamiento del agente}
El ponente mencionó el módulo «Goal-Query»: durante el entrenamiento, el agent se pregunta continuamente «¿qué skill debo completar para poder generalizar después?» y elige objetivos con base en un novelty score. Este proceso cuida el equilibrio entre exploration (novel targets) y exploitation (objetivos conocidos).

\begin{knowledgebox}{Flujo de Goal-Query}
1. Se calcula la novelty de cada stored goal mediante una uncertainty estimate;\\
2. Si la novelty supera el threshold, se programa un exploration goal;\\
3. Con apoyo del skill graph se infieren las dependencias entre goals y se decide si conviene aprender primero los goals básicos.
\end{knowledgebox}

\subsection{Goal space engineering}
Chelsea enfatiza que el goal space debe admitir un escalamiento dinámico, con centro/radio/priority definidos para cada objetivo. Por ejemplo, un complex goal se descompone en primitives; durante el entrenamiento se muestrean varias veces goals de nivel Primitive y luego un scheduler avanza hacia el composite goal.

\begin{tabularx}{\textwidth}{lX}
\toprule
Estrategia & Función \\
\midrule
Primitive sampling & Genera reward-distinct goals a partir de combinaciones de acciones básicas; \\
Composite scheduler & Cuando la primitive success rate alcanza la meta, pondera proporcionalmente la probabilidad de selección del composite goal; \\
Priority annealing & Ajusta dinámicamente la goal priority según la novelty/entropy; \\
\bottomrule
\end{tabularx}

\subsection{Resumen de la sección}
El establecimiento autónomo de objetivos combina el aprendizaje por currículo, Goal-Query y el difficulty scheduler, de modo que el robot puede acumular habilidades de forma gradual sin que una persona le asigne tareas.
\newpage
\section{Exploración segura y operación prolongada}
\subsection{Optimización con restricciones y límites suaves}
\begin{figure}[H]
\centering
\includegraphics[width=0.9\textwidth]{slides-images/slide-06.jpg}
\caption{La Slide 6 trata las restricciones de seguridad como una reward penalty que se usa en el soft policy gradient.}
\end{figure}

\begin{knowledgebox}{Estrategias de exploración segura}
\begin{itemize}
\item Convertir la safety constraint en un penalty term $c(s_t,a_t)$;\\
\item Controlar la intensidad de la penalty con el Lagrangian multiplier $\lambda$;\\
\item Cuando la uncertainty es alta (por ejemplo, ante un new goal), reducir el exploration rate para evitar acciones high-risk.
\end{itemize}
\end{knowledgebox}

\subsection{Monitoreo con un horizon prolongado}
Chelsea enfatiza que la Autonomy no es un único entrenamiento, sino una property de múltiples deployments; por eso se requiere un monitoreo a nivel de sistema:
\begin{itemize}
\item Registrar las trayectorias de success/failure de cada goal;
\item Hacer track de $\Delta reward$ y $\Delta constraint$ para detectar drifting;
\item Ante fallas repetidas, revertir dinámicamente el modelo de reward o reducir el goal space.
\end{itemize}

\begin{warningbox}{Conservar una vía de reversión en el despliegue}
El pipeline de despliegue debe conservar los mecanismos de «revertir a un checkpoint» y «reentrenar la estimación de reward», para poder resolver con rapidez problemas como una hallucination reward o una reset policy que deja de funcionar.
\end{warningbox}

\subsection{Resumen de la sección}
La exploración segura protege el comportamiento mediante una arquitectura de penalty/constraint; la operación prolongada depende de estadísticas de monitoreo y de planes de reversión, y en conjunto conducen a un sistema de Autonomy confiable.
\newpage
\section{Recomendaciones prácticas y operación del sistema}
\subsection{Tablero de monitoreo de Autonomy}
El ponente recomienda construir un tablero con tres indicadores, reward accuracy, reset success y goal coverage, para la revisión diaria.
\begin{tabularx}{\textwidth}{lX}
\toprule
Dimensión & Indicador \\
\midrule
Reward fidelity & classifier precision/recall, clip similarity drift; \\
Reset coverage & tasa de éxito reciente de forward-backward > 80\%; \\
Goal expansion & número de goals activos y goals nuevos por semana; \\
\bottomrule
\end{tabularx}

\subsection{Ciclo de retroalimentación entre etapas}
En cada ejecución de un experimento de Autonomy se registra lo siguiente:
\begin{enumerate}
\item la versión actual del modelo de reward y de la safety penalty;
\item el failure pattern de la Reset policy;
\item la incorporación de nuevos objetivos o habilidades y las decisiones de retirar objetivos anteriores;
\end{enumerate}
Esta información forma un cross-team log que facilita la replicación futura.

\begin{knowledgebox}{Plantilla de Feedback Loop}
1. Pre-scale: confirmar el estado de reward, safety y checkpoint;\\
2. Scale: ejecutar la policy y sincronizar las metrics;\\
3. Post-scale: analizar en retrospectiva el metrics drift y, si se requiere un rollback, activarlo de inmediato.
\end{knowledgebox}

\subsection{Resumen de la sección}
En la práctica, Autonomy requiere un tablero, retroalimentación entre etapas y un plan de reversión para que los resultados de la investigación puedan desplegarse una y otra vez.
\newpage
\section{Exploración multiescala y generalización entre tareas}
\subsection{Estrategia de exploración en varias etapas}
La Slide 07 presenta el multi-stage exploration framework: divide la exploration policy en dos niveles, coarse-grained y fine-grained, que controlan respectivamente los objetivos y las secuencias de acciones de distinta escala. Esta estructura jerárquica permite descubrir macro-level goal cluster y, a la vez, converger automáticamente a los pasos de ejecución micro-level.

\begin{figure}[H]
\centering
\includegraphics[width=0.85\textwidth]{slides-images/slide-07.jpg}
\caption{La Slide 7 explica cómo la multi-stage exploration amplía de forma coordinada el goal space.}
\end{figure}

\begin{knowledgebox}{Multi-Stage Exploration}
\begin{itemize}
\item La coarse policy muestrea goal cluster no vistos;\\
\item La fine policy se encarga de las acciones detalladas dentro de ese cluster;\\
\item Se actualizan periódicamente los cluster centroids based on success rate para mantener la coverage.
\end{itemize}
\end{knowledgebox}

\subsection{Métricas de generalización entre tareas}
La Slide 08 propone una transfer matrix para evaluar la capability de la policy en distintas combinaciones de goal; las métricas principales son success coverage, policy entropy y transfer ratio. El ponente recomienda mostrar en los informes experimentales la performance tanto en los base goals como en los transfer goals, para evitar el sobreajuste a un solo escenario.

\begin{tabularx}{\textwidth}{lX}
\toprule
Métrica & Descripción \\
\midrule
Success coverage & Proporción de stored goals en los que se alcanza el éxito; un valor inferior a 70\% indica una coverage gap; \\
Policy entropy & Supervisa la diversidad de la policy en distintos goals; un valor demasiado bajo indica que es degenerate; \\
Transfer ratio & La improvement que la reward de un goal produce en otro goal; mide la generalización; \\
\bottomrule
\end{tabularx}

\subsection{Resumen de la sección}
La exploración multiescala y las métricas de generalización ayudan al equipo a identificar en qué goals hay que romper cuellos de botella estrechos y ofrecen herramientas para medir la eficacia de la transfer.
\newpage
\section{Integración del sistema y flujo de despliegue}
\subsection{Slide 09-10 Pipeline}
Las Slides 09-10 sintetizan el sistema Autonomy: Reward module, Reset policy, Goal manager y Safety guard trabajan por capas, y los resultados se envían al dashboard para que el operator los revise.

\begin{figure}[H]
\centering
\includegraphics[width=0.9\textwidth]{slides-images/slide-10.jpg}
\caption{La Slide 10 muestra la modular pipeline del sistema Autonomy.}
\end{figure}

\begin{importantbox}{Nodos clave del pipeline de despliegue}
\begin{itemize}
\item Cada Reward model update desencadena un shadow policy test y se hace un sanity check en simulation;\\
\item El Reset controller ejecuta todos los días una recovery routine para asegurar que la forward-backward policy siga siendo confiable;\\
\item El Goal manager actualiza la goal library con base en las live metric y descarta todo goal stale o unsafe;\\
\end{itemize}
\end{importantbox}

\subsection{Retroalimentación entre equipos e intercambio de conocimiento}
El ponente advierte: Autonomy necesita ciclos de retroalimentación entre equipos; por ejemplo, el operator sube cada reset failure case al shared log y, a partir de ello, el equipo de policy ajusta la goal selection o la reward normalization. La documentación y la weekly review meeting ayudan a mantener el conocimiento.

\subsection{Resumen de la sección}
La integración del sistema se centra en organizar el pipeline y en la comunicación entre equipos, para asegurar que Autonomy no solo sea eficaz en la investigación, sino que también responda con rapidez a las anomalías durante el despliegue.
\newpage
\section{Estudio de caso: Vision-guided Autonomy}
\subsection{Descomposición de tareas en escenarios reales}
La Slide 09 muestra un vision-guided manipulation case, que incluye tres módulos: grasp, transport y place; cada módulo, a su vez, anida reward, reset y goal manager autónomos. El ponente enfatiza que los escenarios reales suelen presentar long-tailed failure mode, por lo que es necesario incorporar anomaly detection en el pipeline.

\begin{figure}[H]
\centering
\includegraphics[width=0.85\textwidth]{slides-images/slide-09.jpg}
\caption{La Slide 9 muestra el flujo de trabajo de vision-guided manipulation y enfatiza cómo colaboran los distintos module.}
\end{figure}

\begin{knowledgebox}{División de Vision-guided Autonomy}
\begin{itemize}
\item Grasp module: usa visual affordance para predecir el grasp point;\\
\item Transport module: mueve el objeto mediante una goal-conditioned policy;\\
\item Place module: usa un success classifier para determinar si se completó la colocación en el objetivo.
\end{itemize}
\end{knowledgebox}

\subsection{Observabilidad y respuesta a anomalías}
El ponente sugiere registrar una anomaly metric en cada etapa: confidence drop del modelo de vision, failure rate de la reset policy y goal coverage drop. Al llevar estas métricas a un dashboard, es posible hacer trigger de un automated rollback o de una reward recalibration.

\begin{importantbox}{Matriz de respuesta a anomalías}
\begin{tabularx}{\textwidth}{lX}
\toprule
Tipo de anomalía & Respuesta \\
\midrule
Vision confidence drop & Entrar en safe mode y reducir la exploration rate; \\
Reset failure & Activar de inmediato una manual inspection y revertir al stable checkpoint más reciente; \\
Goal coverage shrink & Ampliar el goal replay buffer y volver a hacer thaw de los failed goal cluster; \\
\bottomrule
\end{tabularx}
\end{importantbox}

\subsection{Resumen de la sección}
Mediante el caso de vision-guided manipulation, se muestran las estrategias de observabilidad y de respuesta a anomalías del Autonomy pipeline en el despliegue fuera de línea y en línea.
\newpage
\section{Síntesis y ampliación}
\subsection{Tabla de síntesis principal}
\begin{tabularx}{\textwidth}{lXl}
\toprule
Dimensión & Idea central & Implicación práctica \\
\midrule
Señal de recompensa & Imágenes, lenguaje y success classifier sustituyen la reward manual & Al diseñar el reward pipeline, incorporar mecanismos de confianza/ensemble; \\
Operación del entorno & Forward-backward controllers y time-aware buffer hacen posible el aprendizaje reset-free & Cada reset policy necesita un registro de recovery y debe vincularse con el scheduler; \\
Gestión de objetivos & Goal-Query y curriculum policy logran una dificultad adaptativa & Mantener un goal repository y registrar novelty/skill dependency; \\
Seguridad y operación & safety penalty + regular monitoring evitan el drift & Construir un dashboard y conservar un canal de rollback; \\
\bottomrule
\end{tabularx}

\subsection{Lecturas adicionales}
\begin{itemize}
\item Sharma et al., ``Autonomous Reinforcement Learning: Formalism and Benchmarking,'' ICLR 2022
\item Eysenbach et al., ``Leave No Trace: Learning to Reset for Safe and Autonomous RL,'' ICLR 2018
\item Pong et al., ``Skew-Fit: State-Covering Self-Supervised RL,'' ICML 2020
\item Nair et al., ``Visual Reinforcement Learning with Imagined Goals,'' NeurIPS 2018
\item Ma et al., ``VIP: Towards Universal Visual Reward and Representation via Value-Implicit Pre-Training,'' ICLR 2023
\end{itemize}

\subsection{Resumen de la sección}
Mediante una exposición sistemática de cinco dimensiones (recompensa, reset, goal, seguridad y monitoreo), esta clase presentó la ruta completa del aprendizaje autónomo de robots y ofreció mecanismos ejecutables de retroalimentación y reversión para una Autonomy de nivel de despliegue.

\end{document}
